% Part: first-order-logic
% Chapter: models-theories
% Section: theories

\documentclass[../../../include/open-logic-section]{subfiles}

\begin{document}

\olfileid{fol}{mat}{the}

\olsection{پہلے درجے دے نظریاں دیاں مثالاں}

\begin{ex}
زبان~$\Lang L_<$ وچ سخت خطی ترتیباں دا نظریہ ہیٹھلے سیٹ
دے مسلّماں نال متعین کیتا جاندا اے:
\begin{align*}
\{\quad & \lforall[x][\lnot x < x], \\
& \lforall[x][\lforall[y][((x < y \lor y <
    x) \lor x = y)]], \\
& \lforall[x][\lforall[y][\lforall[z][((x < y
      \land y < z) \lif x < z)]]] \quad \}
\end{align*}
ایہ ساڈیاں مقصودہ !!{structure}s نوں پوری طرح پکڑدا اے:
ہر سخت خطی ترتیب ایس مسلّمی نظام دا نمونہ اے؛ تے الٹ وی، جے $R$
% PNB-SEM-015: الٹی سمت وچ سخت خطی ترتیب چاہیدی اے، صرف انعکاسی خطی ترتیب نہیں۔
سیٹ $X$ اُتے سخت خطی ترتیب اے، تاں اوہ ساخت $\Struct M$ جیہدے لئی
$\Domain M = X$ تے $\Assign{<}{M} = R$، ایس نظریے دا نمونہ اے۔
\end{ex}

\begin{ex}
گروپاں دا نظریہ، جیہدی زبان وچ $\Obj 1$ (!!{constant}) تے $\cdot$
(دو تھانواں والا !!{function}) نیں، ہیٹھلے مسلّماں نال متعین اے:
\begin{align*}
& \lforall[x][\eq[(x \cdot \Obj 1)][x]]\\
& \lforall[x][\lforall[y][\lforall[z][\eq[(x \cdot (y \cdot z))][((x
          \cdot y) \cdot z)]]]]\\
& \lforall[x][\lexists[y][\eq[(x \cdot y)][\Obj 1]]]
\end{align*}
\end{ex}

\begin{ex}
پیانو حساب دا نظریہ، حساب دی زبان~$\Lang L_A$ وچ ہیٹھلے
جملیاں نال مسلّمی بنیاد پاؤندا اے۔
\begin{align*}
& \lforall[x][\lforall[y][(\eq[x'][y'] \lif \eq[x][y])]]\\
& \lforall[x][\eq/[\Obj 0][x']]\\
& \lforall[x][\eq[(x + \Obj 0)][x]]\\
& \lforall[x][\lforall[y][\eq[(x + y')][(x + y)']]]\\
& \lforall[x][\eq[(x \times \Obj 0)][\Obj 0]]\\
& \lforall[x][\lforall[y][\eq[(x \times y')][((x \times y) + x)]]]\\
& \lforall[x][\lforall[y][(x < y \liff \lexists[z][\eq[(z' + x)][y])]]]\\
\intertext{نال ہی ایس روپ دے سارے جملے وی}
& (!A(\Obj 0) \land \lforall[x][(!A(x) \lif !A(x'))]) \lif \lforall[x][!A(x)]
\end{align*}
چونکہ پچھلے روپ دے بے انتہا جملے نیں، ایس لئی ایہ مسلّمی
نظام لامحدود اے۔ پچھلے روپ نوں \emph{استقرائی خاکہ} آکھدے نیں۔
(اصل وچ، استقرائی خاکہ ایتھے دسے گئے توں تھوڑا ہور پیچیدہ اے۔)

% PNB-SEM-024: ایہ اشارہ استقرائی خاکے توں پہلاں والی آخری حسابی شق ول اے، خود خاکے ول نہیں۔
استقرائی خاکے توں پہلاں والا آخری مسلّمہ~$<$ دی اک \emph{واضح تعریف} اے۔
\end{ex}

\begin{ex}
خالص سیٹاں دا نظریہ ریاضی دیاں بنیاداں وچ، تے ریاضی دے فلسفے
وچ وی، اہم کردار رکھدا اے۔ سیٹ خالص اے جے اوہدے سارے
!!{element}s وی خالص سیٹ ہون۔ ایس لئی خالی سیٹ وی خالص اے،
پر جیہڑے سیٹ دا کوئی !!a{element} خود سیٹ نہ ہووے، اوہ خالص
نہیں۔ سو خالص سیٹ اوہ نیں جیہڑے صرف خالی سیٹ توں بنے ہون
تے انہاں وچ کوئی ”غیر سیٹ بنیادی شے“، یعنی اوہ چیز جیہڑی
آپ سیٹ نہ ہووے، شامل نہ ہووے۔

ہیٹھلے نوں خالص سیٹاں دے نظریے دا مسلّمی نظام سمجھیا جا سکدا اے:
\begin{align*}
& \lexists[x][\lnot \lexists[y][y \in x]]\\
& \lforall[x][\lforall[y][(\lforall[z](z \in x \liff z \in y) \lif 
\eq[x][y])]]\\
& \lforall[x][\lforall[y][\lexists[z][\lforall[u][(u \in z \liff
(\eq[u][x] \lor \eq[u][y]))]]]]\\
& \lforall[x][\lexists[y][\lforall[z][(z \in y \liff \lexists[u][(z \in
        u \land u \in x)])]]]\\
\intertext{نال ہی ایس روپ دے سارے جملے وی} &
\lexists[x][\lforall[y][(y \in x \liff !A(y))]]
\end{align*}
پہلا مسلّمہ آکھدا اے کہ اک ایہو جیہا سیٹ اے جیہدے کوئی
!!{element}s نہیں (یعنی $\emptyset$ موجود اے)؛ دوجا آکھدا اے
کہ سیٹاں دی برابری انہاں دے عنصراں نال طے ہوندی اے؛ تریجا
آکھدا اے کہ ہر دو سیٹاں $X$ تے $Y$ لئی سیٹ $\{X, Y\}$ موجود
اے؛ تے چوتھا کہ ہر سیٹ $X$ لئی سیٹ $\cup X$ موجود اے، جتھے
$\cup X$،~$X$ دے سارے عنصراں دا اتحاد اے۔

آخر وچ دتے گئے !!{sentence}s نوں اکٹھے
\emph{سادہ خاصیت توں سیٹ بنان دا خاکہ} آکھدے نیں۔ ایہ بنیادی
طور تے آکھدا اے کہ ہر $!A(x)$ لئی سیٹ $\Setabs{x}{!A(x)}$
موجود اے؛ پہلی نظر وچ ایہ سچا، مفید، تے شاید ضروری وی لگدا
اے۔ ایہنوں ”سادہ“ اس لئی آکھدے نیں کہ حقیقت وچ ایہ نظریے
نوں پورا ہون جوگا نہیں رہن دیندا: جے $!A(y)$ نوں $\lnot y \in y$
لو، تاں ایہ !!{sentence} ملدا اے:
\[
\lexists[x][\lforall[y][(y \in x \liff \lnot y \in y)]]
\]
تے ایہ !!{sentence} کسے وی !!{structure} وچ پورا نہیں اتردا۔
\end{ex}

\begin{ex}
\emph{حصہ-نظریے} (mereology) وچ \emph{حصہ ہونا} (parthood)
اک بنیادی تعلق اے۔ سیٹاں دے نظریاں وانگ حصے ہون دے وی
نظریے نیں جیہڑے ایس تعلق دیاں وکھ وکھ، تے کدے آپس وچ
ٹکراندیاں، سمجھاں نوں مسلّماں راہیں متعین کردے نیں۔

حصہ-نظریے دی زبان وچ اک ہی دو تھانواں والا محمولی نشان
~$\Obj P$ اے، تے $\Atom{\Obj P}{x, y}$ دا ”مطلب“ ایہ اے
کہ $x$،~$y$ دا حصہ اے۔ ایس تعبیر نوں ذہن وچ رکھ کے، ایس
زبان دی !!a{structure} نوں \emph{حصیاں دی ساخت} آکھدے نیں۔
بے شک اک دو تھانواں والے محمول دی ہر ساخت ایس ناں دی حق دار
نہیں۔ ”حصہ ہون“ نوں پکڑن لئی $\Assign{\Obj P}{M}$ نوں کجھ
شرطاں پوریاں کرنیاں چاہیدیاں نیں، جیہڑیاں اسیں حصہ-نظریے دے
مسلّمے بنا سکدے آں۔ مثال لئی چیزاں وچ حصہ ہون دا تعلق جزوی
ترتیب اے: ہر چیز اپنے آپ دا حصہ اے، بھاویں \emph{غیر سخت}
حصہ؛ دو وکھریاں چیزاں اک دوجے دے حصے نہیں ہو سکدیاں؛
تے کسے چیز دے حصے دا حصہ خود وی اُس چیز دا حصہ اے۔ غور کرو
کہ ایس معنی وچ ”حصہ ہون“ ”ذیلی سیٹ ہون“ نال ملدا اے، پر
”رکن ہون“ نال نہیں، کیوں جے رکن ہون نہ انعکاسی اے نہ متعدی۔
\begin{align*}
& \lforall[x][\Atom{\Obj P}{x,x}] \\
& \lforall[x][\lforall[y][((\Part{x}{y} \land \Part{y}{x})
      \lif \eq[x][y])]] \\
& \lforall[x][\lforall[y][\lforall[z][((\Part{x}{y} \land
        \Part{y}{z}) \lif \Part{x}{z})]]]\\
\intertext{ایس توں ودھ، ہر دو چیزاں لئی حصیاں دا اک مجموعہ ہوندا اے (اک ایہو جیہی چیز جیہدے
  حصہ اوہ دونے ہون، تے ایس لحاظ نال اوہ سبھ توں چھوٹی ہووے)۔}  &
\lforall[x][\lforall[y][\lexists[z][\lforall[u][(\Part{z}{u} \liff
        (\Part{x}{u} \land \Part{y}{u}))]]]]
\end{align*}
ایہ حصہ ہون دے اوہناں بنیادی اصولاں وچوں صرف کجھ نیں جنہاں
اُتے مابعدالطبیعیات دے ماہر غور کردے نیں۔ پر ہور اصول، کجھ
تعریف شدہ تعلق پہلے متعارف کرائے بغیر، جلد ہی بیان یا لکھن
وچ اوکھے ہو جاندے نیں۔ مثال لئی حصہ-نظریے وچ دلچسپی رکھن
والے زیادہ تر ماہر ایہنوں وی معتبر اصول منندے نیں: جدوں وی
کسے چیز~$x$ دا اک سخت حصہ~$y$ ہووے، تاں اوہدا اک ہور
حصہ~$z$ وی ہووے جیہدا~$y$ نال کوئی مشترک حصہ نہ ہووے،
تے $y$ تے $z$ دے مجموعے نال~$x$ بنے۔
\end{ex}

\end{document}
